\textbf{T1: Sunny (10 pts)}

You are asked to study the features of the brightly lit circle and dark rings in
the figures below. Make your calculations for an idealized situation: the chair
leg is strictly cylindrical of radius $a$, strictly vertical, with a perfectly
smooth, cylindrical, and perfectly reflecting surface. You may make any
additional model assumptions and approximations you deem reasonable that will
simplify your calculations.

a) (5 pts) Determine how the illuminance surplus $I(r,\theta)$ inside the
brightly lit circle on the floor depends on the polar coordinates $r \gg a$ and
$\theta$. The illuminance quantifies the amount of incoming light per area. By
``surplus'' we mean the additional illuminance introduced due to the presence of
the cylinder. Express the answer in terms of $I_0$ defined as the illuminance
difference between points A and B in the figure.

\begin{center}\includegraphics[width=0.49\linewidth]{figures/EuPhO_2025_Q1_fig1.jpg}\end{center}

b) (5 pts) In the following figure some fingers are blocking some of the light
from reaching the chair leg. Let $R(\theta)$ denote the radial distance of the
middle dark ring as a function of the angle $\theta$ and let $R_\mathrm{min}$ be
the minimal value of $R(\theta)$. Determine $R(\theta) - R_\mathrm{min}$.

\begin{center}\includegraphics[width=0.49\linewidth]{figures/EuPhO_2025_Q1_fig2.jpg}\end{center}
